\chapter{Terzo capitolo}
\label{cap:tre}

    \section{Prima sezione}
    \label{sec:tre_uno}

        \lipsum[22][1-5] La formulazione adottata segue quella di~\cite{ferrari2019rassegna}.

        \begin{equation}
            \label{eq:esempio_uno}
            \bar{x} = \frac{1}{n} \sum_{i=1}^{n} x_i ,
        \end{equation}

        dove $n$ indica la numerosità del campione. L'Equazione~\ref{eq:esempio_uno} \lipsum[23][1-2]

        \lipsum[24][1-3] Gli elementi considerati per la \ac{CCC} sono i seguenti:

        \begin{itemize}
            \item \textbf{Primo elemento:} \lipsum[25][1]
            \item \textbf{Secondo elemento:} \lipsum[26][1]
            \item \textbf{Terzo elemento:} \lipsum[27][1]
        \end{itemize}

    \section{Seconda sezione}
    \label{sec:tre_due}

        \lipsum[28][1-4] Le fasi della procedura sono riassunte di seguito~\cite{rossi2020metodo}:

        \begin{enumerate}[label=(\alph*)]
            \item \lipsum[29][1]
            \item \lipsum[30][1]
            \item \lipsum[31][1]
        \end{enumerate}

        \lipsum[32][1-3] Il valore riportato è pari a \SI{12.5}{\percent} del totale.
